\documentclass{article}

\usepackage[a4paper, left=0.855in, right=0.855in, top=0.75in, bottom=0.75in]{geometry}
\usepackage{amsmath}
\usepackage{amssymb}
\usepackage{fontspec}
\usepackage{titling}
\usepackage{titlesec}
\usepackage{indentfirst} 
\usepackage[style=authoryear, backend=biber]{biblatex}
\addbibresource{ds_a2_q4.bib}  % 指定 .bib 文件

\title{DS\_A2\_Q4}
\author{Dongju Ma}
\date{Trimester 1, 2025}

% 设置段落缩进
\setlength{\parindent}{2em} % 设置缩进量为2字符

% 设置section格式
\titleformat{\section}[hang]
    {\normalfont\Large\bfseries}{\thesection}{1em}{}
\titlespacing*{\section}{0pt}{3.5ex plus 1ex minus .2ex}{2.3ex plus .2ex}

\linespread{1.25}

\begin{document}

\maketitle

\section{Instructions of PCR}
Principal component regression (PCR) is a regression technique that is based on principal component analysis (PCA). 
The main idea of PCR is to reduce the dimensionality of the data by projecting it onto a lower-dimensional subspace. 
It can be used in processing multiple predictors which are quantitative continous variables, especifically useful when the number of predictors is large compared to the number of observations, as it can help to reduce the risk of overfitting.
PCR is a kind of regression model, with two steps: 
\begin{itemize}
    \item Perform PCA on the predictors to obtain the principal components.
    \item Use the principal components as predictors in a linear regression model.
\end{itemize}


\section{How PCR exactly works}

\subsection{Step 1: Perform PCA}
As we mentioned above, we should perform PCA on the predictors, so what is PCA?
Principal Component Analysis (PCA) is a technique that transforms correlated predictors into uncorrelated principal components through eigendecomposition of the covariance matrix, where the eigenvectors define the transformation axes and the eigenvalues indicate their importance 

For example if we have \(p\) predictors \(X_1, X_2, \dots, X_p\) belong to a dimension space \(R^p\), we can transform them into \(m\) principal components \(Z_1, Z_2, \dots, Z_M, (M \leq p)\) belong to the same dimension space \(R^p\) via linear combinations:

\begin{equation}
    Z_m = \sum_{j=1}^{p} \phi_{jm} {X}_j, \quad m = 1, 2, \dots, M, \quad (M \leq p)
\end{equation}

In Formula (1), \(Z_m\) is the \(m\)th principal component, \(X_j\) is the \(j\)th predictor, and \(\phi_{jm}\) is the \(j\)th loading of the \(m\)th principal component, which is the \(j\)th element of the \(m\)th eigenvector of the covariance matrix of the predictors. 
To determine \(\phi_{jm}\), we should calculate the covariance matrix of the predictors, 
\begin{equation}
    Cov(X) = \frac{1}{n - 1} \sum_{i=1}^{n} (X_i - \bar{X}) (X_i - \bar{X})^T
\end{equation}
We typically use the unbiased estimator to calculate the covariance matrix, so the denominator is \(n - 1\). 

Then we should decompose the covariance matrix to obtain eigenvectors and eigenvalues, 

\begin{equation}
    Cov(X) = P \Lambda P^T
\end{equation}
where \(P\) is a matrix whose columns are the eigenvectors of the covariance matrix, and \(\Lambda\) is a diagonal matrix whose diagonal elements are the eigenvalues of the covariance matrix. 
\(P\) can be loaded onto the original observations, which is considerd that \(\phi_{jm}\) represents the importance of the \(j\)th predictor in the \(m\)th principal component.

Besides, standardisation of predictors in pre-processing is recommended as PCA is scale-sensitive. 

\subsection{Step 2: Linear regression with the principal components}
As \(Z_m\) can be considered as the score of principal components based on subsections of original observations and uncorrelated, we can use them as predictors in a linear regression model.
\begin{equation}
    Y_i = \theta_0 + \sum_{m=1}^{M} \theta_m Z_im + \epsilon_i, \quad i = 1, 2, \dots, n
\end{equation}
where \(Y_i\) is the response variable, \(\theta_0\) is the intercept, \(\theta_m\) is the coefficient of the \(m\)th principal component, \(Z_im\) is the \(m\)th principal component of the \(i\)th observation, and \(\epsilon_i\) is the error term of the \(i\)th observation. 

This model means that now we just consider the principal components as predictors and how they could affect the response variable by linear regression.

\subsection{How to choose the number of principal components}
The number of principal components \(M\) is a hyperparameter that we need to choose before performing PCR to the original datasets. 
According to \parencite{hastie_2009_the}, \(M\) is determined by the eigenvalues of the covariance matrix, which means PCR would discard the principal components with small eigenvalues.
There are several optional methods to determine the number of principal components, such as cross-validation, AIC and BIC, which could be chosen as needed.

\section{Scope and Limitation}
As the regression model performs, now we solve the multiple correlation problem by using PCR, but that doesn't mean we can avoid overfitting or PCR is always better than other regression models.

\subsection{Advantage of PCR}
\begin{itemize}
    \item Reduce the dimensionality of the data, which can help to avoid overfitting.
    \item Solve the multicollinearity problem, as the principal components are uncorrelated.
    \item Can be used in processing multiple predictors which are quantitative continous variables.
\end{itemize}
\subsection{Disadvantage of PCR}
\begin{itemize}
    \item The principal components are linear combinations of the original predictors, which may be hard to interpret.
    \item The number of principal components should be determined by the eigenvalues of the covariance matrix, which may be subjective. Other methods like SVD could be used to solve this problem.
    \item The principal components are uncorrelated, but the original predictors may have some interactions.
\end{itemize}

\printbibliography[title={References}]

\end{document}